\section{The invisible creation in the history of salvation}
\label{sec:scripture}

The angels belong to God's creation, and their greatness is a received
greatness. Scripture places them before the divine throne, beside a traveler,
at a prison door, and at the empty tomb. In each setting their action directs
attention beyond themselves. The question in Hebrews joins their spiritual
nature to their mission: ``Are they not all ministering spirits, sent to
minister for them who shall receive the inheritance of salvation?''\footnote{Hebrews
1:14, Douay--Rheims, Challoner revision, Project Gutenberg 1581. Biblical
quotations in this study follow this witness unless stated otherwise.}
To understand an angel is therefore to understand a creature who receives
being from God and can become a servant of God's gifts to others.

\subsection{The messenger and the One who sends}

At Mambre, Abraham receives three visitors; the narrative names the Lord's
appearance, describes the visitors as men, and subsequently identifies the
two who enter Sodom as angels. The story moves between the divine promise
and its creaturely bearers. Abraham's hospitality and Sarah's promised child
remain its center. The mysterious manner of the visitation raises a question
that Aquinas will address precisely: how can a spirit appear and act through
a body without becoming the body's soul?\footnote{Genesis 18:1--15;
19:1--3; \ST{I q.~51 aa.~2--3}. The philosophical question is Aquinas's
treatment of the scriptural appearance.}

Jacob's ladder likewise joins angelic movement with God's initiative.
The angels ascend and descend, but the promise of land, descendants, blessing,
and protection comes from the Lord. Heaven is open because God makes himself
known. The ladder does not describe a technique by which Jacob climbs beyond
his creaturehood; he is asleep when he receives the vision.\footnote{Genesis
28:10--17.}

In Exodus, the angel promised to Israel goes before the people and bears
God's name. The command to heed him depends on the sender's authority.
This pattern helps explain why a messenger can speak with great authority
without possessing a second source of divine sovereignty. It also prepares
a distinction essential to Christian worship: the honor given a servant on
account of God and the adoration owed to God have different objects.\footnote{Exodus
23:20--23; compare Apocalypse 19:10; 22:8--9.}

\subsection{Fire, wings, and the holy God}

Isaias sees the Lord enthroned and the seraphim attending him. Their six
wings conceal as well as carry; their proclamation fills the temple; one
brings the burning coal that touches the prophet's lips. The scene joins
adoration, purification, and mission. Isaias's readiness to be sent follows
the cleansing of his unclean lips. The seraph's action serves that divine
initiative.\footnote{Isaias 6:1--8.}

The image gives Christian reflection more than an inventory of celestial
features. Fire can signify the communication of holiness; wings can signify
readiness and elevation; the covered face can signify reverence before what
exceeds a creature's grasp. Dionysius will develop these meanings carefully.
His interpretation does not reduce the vision to a diagram of anatomy.
It asks how bodily images teach embodied readers about intellectual
creatures.\footnote{\emph{Celestial Hierarchy} II.1--5; XIII; XV.2--4.}

The psalms place angels within the universal summons to praise. Psalm 148
addresses the heavenly hosts together with sun, moon, and stars, and then
names the reason for their existence: God's command. Psalm 102 [Hebrew 103]
joins angelic strength to obedient attention. Their power is intelligible
through its orientation to the Lord's word.\footnote{Psalms 148:1--5;
102:20--22 [Hebrew 103:20--22].}

\subsection{Personal help and the government of peoples}

Raphael's disclosure at the end of Tobias gathers the journey into a single
account of providence. He has presented prayer, accompanied human action,
and brought healing by God's command. His departure returns the family to
praise of the sender. The helper was real even when his identity was hidden;
the help remains a gift even after his identity is disclosed.\footnote{Tobias
12:12--22. Vulgate/Douay numbering is used; translations following a different
Greek recension can differ in both wording and verse distribution.}

Daniel sets angelic ministry against a wider horizon. The heavenly court
contains an immense multitude; Michael is described as a prince who stands
for Daniel's people; the vision joins earthly conflict with a struggle beyond
Daniel's sight. The scale changes from one household to peoples and kingdoms,
but the subject remains God's rule and his servants' participation in it.\footnote{Daniel
7:9--10; 10:12--21; 12:1--3.}

Christ's warning against despising the little ones makes the same order
personally demanding. Their angels behold the Father's face. The dignity
of the little one is thus joined to the heavenly life; angelology has a
moral consequence for how a neighbor is treated. The vision of God and
care for a human being appear together in one saying, a conjunction Aquinas
will preserve in his account of mission.\footnote{Matthew 18:10;
\ST{I q.~112 a.~1 ad~3; q.~113 aa.~1--2}.}

\subsection{The Son and his servants}

Gabriel announces the births of John and Jesus; the heavenly host praises
God at Bethlehem; the messengers at the tomb announce the Resurrection.
They illuminate events whose saving source is Christ. Gabriel communicates
the promise, while the Holy Spirit's power brings about the virginal
conception. This difference between announcing and accomplishing prevents
the messenger's brilliance from obscuring the mystery he announces.\footnote{Luke
1:19--38; 2:8--15; 24:1--8.}

Hebrews opens by distinguishing the Son from the angels who adore him.
Colossians places the invisible powers among the things created through
him. Even the names that sound most sovereign---thrones, dominations,
principalities, powers---name beings whose existence depends on Christ.
They cannot be a route past him to a more ultimate divine reality.\footnote{Hebrews
1:1--14; Colossians 1:15--20; 2:18--19.}

The Apocalypse gives Michael's victory its Christological setting: the
accuser is overcome, and the faithful conquer through the Lamb's blood.
At the close of the book the angel refuses John's adoration and identifies
himself as a fellow servant: ``Adore God.'' The heavenly world is full of
real agents, yet the worship toward which they direct the Church is
undivided.\footnote{Apocalypse 12:7--12; 22:8--9.}

Human destiny is communion with this heavenly company, not a change of human
nature into angelic nature. Jesus compares the risen to angels concerning
marriage; Hebrews names angels and perfected human spirits distinctly in
the one heavenly assembly. The Christian hope retains bodily resurrection
while opening human life to communion beyond the visible world.\footnote{Matthew
22:29--32; Hebrews 12:22--24.}
